\documentclass{book}
\usepackage[english]{babel}
\usepackage[utf8]{inputenc}
\usepackage{geometry,amsmath,amssymb,hyperref,xcolor,latexsym,pslatex}
\geometry{portrait}

%%%%%%%%%% Start TeXmacs macros
\newcommand{\assign}{:=}
\newcommand{\cdummy}{\cdot}
\newcommand{\mathd}{\mathrm{d}}
\newcommand{\mathi}{\mathrm{i}}
\newcommand{\mathpi}{\pi}
\newcommand{\tmcolor}[2]{{\color{#1}{#2}}}
\newcommand{\tmem}[1]{{\em #1\/}}
\newcommand{\tmop}[1]{\ensuremath{\operatorname{#1}}}
\newcommand{\tmrsup}[1]{\textsuperscript{#1}}
\newcommand{\tmstrong}[1]{\textbf{#1}}
\newenvironment{proof}{\noindent\textbf{Proof\ }}{\hspace*{\fill}$\Box$\medskip}
\newtheorem{axiom}{Axiom}
\newtheorem{corollary}{Corollary}
\newtheorem{lemma}{Lemma}
%%%%%%%%%% End TeXmacs macros

%


\begin{document}

\title{Build Quantum Mechanics from Scratch}

\date{June 8, 2026}

\maketitle

\section{Introduction}

We derive the Schrödinger's equation of quantum mechanics from several basic
experimental facts. These facts are treated as axioms. In fact, what we will
get is a generalized Schrödinger equation, which is the most generic equation
that these axioms can imply.

\subsection{Conventions}

\begin{itemize}
  \item {\tmstrong{Definitions}} are in bold font.
  
  \item {\tmem{Important}} statements are in italic font.
  
  \item Only important equations are numbered.
  
  \item \tmcolor{red}{Questions} are in red color.
  
  \item {\color[HTML]{008000}Conclusions} are in green color.
\end{itemize}

\subsection{Abbreviations}

Because we are doing high-dimensional calculus, abbreviations are essential
for not falling in the ``debauchery of formulae''.

When $x, y \in \mathbb{R}^d$, denote $xy = x \cdot y = \sum_{\alpha}
x^{\alpha} y^{\alpha}$. Thus, $x^2 = x \cdot x = \sum_{\alpha}
(x^{\alpha})^2$. Furthermore, we use multi-dimensional index $\alpha =
(\alpha_1, \ldots, \alpha_n)$ for abbreviating partial derivative
$(\partial_{\alpha_1} \cdots \partial_{\alpha_n} f) (x)$ to
$\partial_{\alpha}^n f (x)$, and abbreviating tensor (generalization of vector
to multiple indices) $A^{\alpha_1 \cdots \alpha_n} \tmop{toA}^{\alpha}$, and
$B_{\beta_1 \cdots \beta_n}$ to $B_{\beta}$. For example, the product
$v^{\alpha_1} \times \cdots \times v^{\alpha_n}$ can be abbreviated as
$v^{\alpha}$.

Einstein's convention of summation is an elegant abbreviation for
high-dimensional calculus and tensor analysis. It neglects the
$\Sigma_{\alpha}$ in $\sum_{\alpha} v_{\alpha} w^{\alpha}$, abbreviated to be
$v_{\alpha} w^{\alpha}$. The indices are ``balanced'', one superscript, and
the other subscript. The convention does not apply to the ``unbalanced''
summation such as $\sum_{\alpha} v^{\alpha} w^{\alpha}$, because both indices
are superscripts.

\section{Probability Interpretation and Wavefunction}

In classical Newtonian world, the states of a physical system with $n$
particles are represented by elements in $\mathbb{R}^{6 n}$, or
{\tmstrong{phase}}s, $3 n$ dimensions for positions and $3 n$ for velocities
(or momenta). The time evolution is thus a trajectory in the phase space,
namely a map $\mathbb{R} \rightarrow \mathbb{R}^{6 n}$ (with time as the
domain).

But for a quantum system, the states are represented by wavefunctions. A
{\tmstrong{wavefunction}} is a map from the positions of particles, together
with time, to complex plane. Thus, a wavefunction maps $\mathbb{R}^d \times
\mathbb{R} \rightarrow \mathbb{C}$ where $d = 3 n$. The first axiom claims how
wavefunctions characterize a quantum system.

\begin{axiom}[Probability Interpretation]
  \label{axiom:prob}The states of a quantum system are represented by
  wavefunctions. And given a wavefunction $\varphi$, the probabilistic density
  that the particles are found in positions $x$ at time $t$ is given by $|
  \varphi (x, t) |^2 = \varphi^{\ast} (x, t) \varphi (x, t)$.
\end{axiom}

\section{Superposition Principle and Time Evolution}\label{section:Superposition Principle and Time Evolution}

The second axiom of quantum mechanics, superposition principle, claims that
operations on wavefunctions shall be linear (so that wavefunctions can be
superpositioned).

\begin{axiom}[Superposition Principle]
  \label{axiom:sup}Physical laws that operate on quantum states shall be
  linear.
\end{axiom}

An implication of superposition principle is how quantum states (precisely,
their wavefunctions) evolve with time. Axiom \ref{axiom:sup} implies that the
equation of time evolution (as a physical law that operates on a quantum
state) shall be linear: $\partial \varphi / \partial t = L (\varphi)$ where
the operation $L$ is linear.\footnote{You may wonder why it is not
$(\partial^2 \varphi / \partial t^2) = L (\varphi)$ instead, which is linear
in wavefunction too. If so, wavefunction at a given time cannot fully
characterize the system at that time, just like knowing the particle position
at time $t$ is insufficient for predicting the subsequent position at $(t +
\mathd t)$ in classical physics, for which velocity or momentum is also
essential. But just like the time evolution of a classical system at a given
time is fully characterized by the phase at that time, so that the
(Hamiltonian) dynamics is first order in time derivative, wavefunction has
fully characterized a quantum system (claimed in axiom \ref{axiom:prob}) so
that its time evolution is first order in time derivative too.}
Mathematically, linearity imports a kernel $r : \mathbb{R}^d \times
\mathbb{R}^d \rightarrow \mathbb{C}$ such that
\begin{equation}
  \mathi \frac{\partial \varphi}{\partial t}  (x, t) = \int_{\mathbb{R}^d}
  \mathd y {\color[HTML]{000000}r (x, y, t)} \varphi (y, t) .
  \label{equ:superposition}
\end{equation}
The right hand side can be seen as a generalization of linear transformation
in $\mathbb{R}^n$ like vector-matrix product $\sum_j r_{ij} \varphi_j$. The
imaginary $\mathi$ is employed for convenience.

A direct result of probability interpretation (axiom \ref{axiom:prob}) is that
probabilistic density shall be normalized. Namely, for any wave-function
$\varphi$ and any $t \in \mathbb{R}$, we shall have
\begin{equation}
  \int_{\mathbb{R}^d} \mathd x \varphi^{\ast} (x, t) \varphi (x, t) = 1.
  \label{eq:probtoself}
\end{equation}
Taking derivative on $t$ gives
\[ \int_{\mathbb{R}^d} \mathd x \frac{\partial \varphi^{\ast}}{\partial t} 
   (x, t) \varphi (x, t) + \int_{\mathbb{R}^d} \mathd x \varphi^{\ast} (x, t)
   \frac{\partial \varphi}{\partial t}  (x, t) = 0. \]
Replacing $(\partial \varphi / \partial t)$ by equation
\ref{equ:superposition} (and its complex conjugation),
\[ \int_{\mathbb{R}^d} \mathd x \int_{\mathbb{R}^d} \mathd y \varphi (x, t) 
   [{\color[HTML]{000000}r^{\ast} (x, y, t)} - {\color[HTML]{000000}r (y, x,
   t)}] \varphi^{\ast} (y, t) = 0. \]
Since $\varphi$ is arbitrary, we obtain
\begin{equation}
  {\color[HTML]{000000}r^{\ast}  (x, y, t)} = {\color[HTML]{000000}r (y, x,
  t)} .
\end{equation}
That is, complex conjugating $r$ is simply swapping its arguments. We call
such function {\tmstrong{Hermitian}}. The two arguments of $r$ are not
independent.

Probability interpretation (axiom \ref{axiom:prob}), together with
superposition principle (axiom \ref{axiom:sup}), is the direct result of the
double-slit experiment of electron. Details can be found in Feynman's Lectures
on Physics, Vol 3, chapter 1.

\section{Wavefunction Is Rapidly Decreasing}\label{section:Wavefunction Is Rapidly Decreasing}

Before calculations, we first address where wavefunctions live. Probability
interpretation (axiom \ref{axiom:prob}) demands that wavefunctions are
square-integrable. Namely, wavefunction is in the square-integrable space $L^2
(\mathbb{R}^d)$. But this is far from sufficient. Many mathematical tools are
essential for evaluating quantum mechanics. Two of the many are Fourier
transform and Taylor series.

It is
\href{https://en.wikipedia.org/wiki/Fourier\\\\\\\\\\\\\\\\\\\\\\\\\\\\\\\_transform}{well
known}, however, that the Fourier transform of a square-integrable function
may not be square-integrable again, so that its inverse Fourier transform may
not exist. Since Fourier transform is basic in quantum mechanics, we shall
seek for a smaller space in which wavefunctions live. An ideal substitution is
the (complex) rapidly decreasing functions. {\tmstrong{Rapidly decreasing
function}} is a smooth function $f : \mathbb{R}^d \rightarrow \mathbb{C}$ that
decays ``exponentially fast'' at infinity. Precisely, for any $m$-order
polynomial $P_m$ and any $m$-order partial derivative $\partial^n$,
\footnote{For example, $\partial_{\alpha} \partial_{\beta}^2$ is $3$-order
partial derivative, and $\partial_{\alpha} \partial_{\beta}^2
\partial_{\gamma}^9$ is $12$-order.} with integers $m, n \geqslant 0$, we have
\[ \lim_{\|x\| \rightarrow \infty} |P_m (x) \partial^n f (x) | = 0. \]
Fourier transform on a rapidly decreasing function results in a rapidly
decreasing function.\footnote{This was first found by Laurent Schwartz in
1949, in his paper ``Théorie des distributions et transformation de
Fourier''. The space of rapidly decreasing functions is now called
{\tmstrong{Schwartz space}}, denoted by $\mathcal{S} (\mathbb{R}^d)$. Thus,
Fourier transform is a linear automorphism in Schwartz space.} Denote
$\mathcal{S} (\mathbb{R}^d)$ as the collection of all the rapidly decreasing
functions that map from $\mathbb{R}^d$ to $\mathbb{C}$.

Rapidly decreasing functions are dense in square-integrable space, meaning
that for any $f \in L^2 (\mathbb{R}^d)$ and any $\varepsilon > 0$, there is a
$g \in \mathcal{S} (\mathbb{R}^d)$ such that
\[ \int_{\mathbb{R}^d} \mathd x |f (x) - g (x) |^2 < \varepsilon . \]
For example, for any wavefunction $f \in L^2 (\mathbb{R}^d)$, when we measure
the probability on any area of positions $U \subset \mathbb{R}^d$, we can use
its approximation $g \in \mathcal{S} (\mathbb{R}^d)$ instead, because the
difference is bounded by
\[ \left| \int_U \mathd x|f (x) |^2 - \int_U \mathd x|g (x) |^2 \right|
   \leqslant \int_U \mathd x |f (x) - g (x) |^2 \leqslant \int_{\mathbb{R}^d}
   \mathd x |f (x) - g (x) |^2 < \varepsilon . \]
The first inequality is equivalent to $|\|f\|-\|g\|| \leqslant \|f - g\|$,
where the norm is defined as $\|f\| \assign \int_U \mathd x |f (x) |^2$,
recognized as the $L^2$-norm on $U$. It states that the difference between the
two sides of a triangle is less than that of the third side. It is in this
sense that the substitution $L^2 (\mathbb{R}^d) \rightarrow \mathcal{S}
(\mathbb{R}^d)$ is plausible. Proof of the statement that rapidly decreasing
entire functions are dense in square-integrable space is given in appendices
\ref{appendix:dense}.

This question is answered in appendix \ref{appendix:schwartz-preserve}.

\section{Path Integral Formalism}\label{section:Path Integral Formalism}

We are trying to derive a generic path integral formalism. Given a small
$\Delta t > 0$, time evolution (equation \ref{equ:superposition}) gives
\[ \varphi (x, t + \Delta t) = \int_{\mathbb{R}^d} \mathd y [\delta (x - y) -
   \mathi {\color[HTML]{000000}r (x, y, t)} \Delta t] \varphi (y, t) +
   \omicron (\Delta t) . \]
We are to convert the $[\cdots]$ part into exponential. To do so, we take the
inverse Fourier transform
\[ \delta (x - y) = \int_{\mathbb{R}^d} \frac{\mathd k}{(2 \mathpi)^d} \exp
   (\mathi k (x - y)), \]
and
\begin{equation}
  {\color[HTML]{000000}r (x, y, t)} = \int_{\mathbb{R}^d} \frac{\mathd k}{(2
  \mathpi)^d} \exp (\mathi k (x - y)) {\color[HTML]{000000}\hat{r}  (k, y,
  t)}, \label{eq:r-fourier}
\end{equation}
in which

\footnote{Indeed, by inserting equation \ref{eq:hamiltonian} into equation
\ref{eq:r-fourier}, we get
\[ {\color[HTML]{000000}r (x, y, t)} = \int_{\mathbb{R}^d} \frac{\mathd k}{(2
   \mathpi)^d} \exp (\mathi k (x - y)) \times \left[ \int_{\mathbb{R}^d}
   \mathd x' \exp (- \mathi k (x' - y)) {\color[HTML]{000000}r (x', y, t)}
   \right] . \]
Re-arrange the right hand side as
\[ \int_{\mathbb{R}^d} \mathd x' {\color[HTML]{000000}r (x', y, t)}
   \int_{\mathbb{R}^d} \frac{\mathd k}{(2 \mathpi)^d} \exp (\mathi k (x - x'))
   = \int_{\mathbb{R}^d} \mathd x' {\color[HTML]{000000}r (x', y, t)} \delta
   (x - x'), \]
which goes back to ${\color[HTML]{000000}r (x, y, t)}$, indicating that
equations \ref{eq:r-fourier} and \ref{eq:hamiltonian} are consistent.
\tmcolor{red}{But, before Fourier transform, we have to prove that $r$ can be
approximated by rapidly decreasing functions.}}

\

\footnote{Alternatively, we can define
\begin{equation}
  {\color[HTML]{000000}\check{r}  (x, k, t)} \assign \int_{\mathbb{R}^d}
  \mathd y \exp (- \mathi k (x - y)) {\color[HTML]{000000}r (x, y, t)} .
  \label{eq:r-fourier-alt}
\end{equation}
Thus,
\[ {\color[HTML]{000000}r (x, y, t)} = \int_{\mathbb{R}^d} \frac{\mathd k}{(2
   \mathpi)^d} \exp (\mathi k (x - y)) {\color[HTML]{000000}\check{r}  (x, k,
   t)} . \]
Indeed, plugging ${\color[HTML]{000000}\check{r}  (x, k, t)}$ into the right
hand side of ${\color[HTML]{000000}r (x, y, t)}$,
\[ {\color[HTML]{000000}r (x, y, t)} = \int_{\mathbb{R}^d} \frac{\mathd k}{(2
   \mathpi)^d} \exp (\mathi k (x - y)) \times \left[ \int_{\mathbb{R}^d}
   \mathd y' \exp (- \mathi k (x - y')) {\color[HTML]{000000}r (x, y', t)}
   \right] . \]
Re-arrange the right hand side as
\[ \int_{\mathbb{R}^d} \mathd y' {\color[HTML]{000000}r (x, y', t)}
   \int_{\mathbb{R}^d} \frac{\mathd k}{(2 \mathpi)^d} \exp (\mathi k (y' - y))
   = \int_{\mathbb{R}^d} \mathd y' {\color[HTML]{000000}r (x, y', t)} \delta
   (y' - y), \]
which goes back to ${\color[HTML]{000000}r (x, y, t)}$ again. The $\hat{r}$
and $\check{r}$ are the Fourier transform of $r$ performed on each of its
arguments respectively. In fact, $\hat{r}$ and ř are the same object. Indeed,
recalling the Hermitianity of $r$, we have
\[ {\color[HTML]{000000}\hat{r}^{\ast}  (k, y, t)} = \int_{\mathbb{R}^d}
   \mathd x \exp (\mathi k (x - y)) {\color[HTML]{000000}r^{\ast}  (x, y, t)}
   = \int_{\mathbb{R}^d} \mathd x \exp (\mathi k (x - y))
   {\color[HTML]{000000}r (y, x, t)} . \]
Exchanging $x$ and $y$ makes
\[ {\color[HTML]{000000}\hat{r}^{\ast}  (k, x, t)} = \int_{\mathbb{R}^d}
   \mathd y \exp (- \mathi k (x - y)) {\color[HTML]{000000}r (x, y, t)}, \]
which is just the ${\color[HTML]{000000}\check{r}  (x, k, t)}$. So, we have
\begin{equation}
  {\color[HTML]{000000}\hat{r}^{\ast}  (k, x, t)} =
  {\color[HTML]{000000}\check{r}  (x, k, t)} .
\end{equation}
Once again, we find that the two arguments of ${\color[HTML]{000000}r (x, y,
t)}$ are not independent.}
\begin{equation}
  {\color[HTML]{000000}\hat{r}  (k, y, t)} \assign \int_{\mathbb{R}^d} \mathd
  x \exp (- \mathi k (x - y)) {\color[HTML]{000000}r (x, y, t)} .
  \label{eq:hamiltonian}
\end{equation}
Then, the $[\cdots]$ part is converted into exponential by

\begin{align}
  {}[\cdots] = & \int_{\mathbb{R}^d} \frac{\mathd k}{(2 \mathpi)^d} \exp
  (\mathi k (x - y))  [1 - \mathi {\color[HTML]{000000}\hat{r} (k, y, t)}
  \Delta t] \nonumber\\
  = & \int_{\mathbb{R}^d} \frac{\mathd k}{(2 \mathpi)^d} \exp \{\mathi k (x -
  y) - \mathi {\color[HTML]{000000}\hat{r} (k, y, t)} \Delta t\} + \omicron
  (\Delta t) \nonumber
\end{align}

Plugging back to $\varphi (x, t + \Delta t)$results in
\[ \varphi (x, t + \Delta t) = \int_{\mathbb{R}^d} \mathd y
   \int_{\mathbb{R}^d} \frac{\mathd k}{(2 \mathpi)^d} \exp \{\mathi k (x - y)
   - \mathi {\color[HTML]{000000}\hat{r} (k, y, t)} \Delta t\} \varphi (y, t)
   + \omicron (\Delta t) . \]
Now, we have converted the $[\cdots]$ part into exponential, as a starting
point of constructing path integral.

After re-denoting $x_1 \assign x$, $x_0 \assign y$, $k_0 \assign k$, it
becomes
\[ \varphi (x_1, t + \Delta t) = \int_{\mathbb{R}^d} \mathd x_0 
   \int_{\mathbb{R}^d} \frac{\mathd k_0}{(2 \mathpi)^d} \exp \{\mathi k_0 (x_1
   - x_0) - \mathi {\color[HTML]{000000}\hat{r} (k_0, x_0, t)} \Delta t\}
   \varphi (x_0, t) + \omicron (\Delta t) . \]
The same (replacing $x_1$ by $x_2$, and $x_0$ by $x_1$),
\[ \varphi (x_2, t + 2 \Delta t) = \int_{\mathbb{R}^d} \mathd x_1 
   \int_{\mathbb{R}^d} \frac{\mathd k_1}{(2 \mathpi)^d} \exp \{\mathi k_1 (x_2
   - x_1) - \mathi {\color[HTML]{000000}\hat{r} (k_1, x_1, t)} \Delta t\}
   \varphi (x_1, t + \Delta t) + \omicron (\Delta t) . \]
By inserting $\varphi (x_1, t + \Delta t)$, we find

\begin{align}
  \varphi (x_2, t + 2 \Delta t) = & \int_{\mathbb{R}^d} \mathd x_1 
  \int_{\mathbb{R}^d} \frac{\mathd k_1}{(2 \mathpi)^d} \exp \{\mathi k_1 (x_2
  - x_1) - \mathi {\color[HTML]{000000}\hat{r} (k_1, x_1, t)} \Delta t\}
  \times \nonumber\\
  \times & \int_{\mathbb{R}^d} \mathd x_0  \int_{\mathbb{R}^d} \frac{\mathd
  k_0}{(2 \mathpi)^d} \exp \{\mathi k_0 (x_1 - x_0) - \mathi
  {\color[HTML]{000000}\hat{r} (k_0, x_0, t)} \Delta t\} \varphi (x_0, t)
  \nonumber\\
  & + \omicron (\Delta t) . \nonumber
\end{align}

After repeating this $N$ times, we arrive at
\begin{equation}
  \varphi (x_N, t + N \Delta t) = \int D (k, x) \exp (\mathi S (k, x)) \varphi
  (x_0, t) + \omicron (\Delta t), \label{eq:pathint}
\end{equation}
in which
\begin{equation}
  \int D (k, x) \assign \prod_{i = 0}^{N - 1} \int_{\mathbb{R}^d} \frac{\mathd
  k_i}{(2 \mathpi)^d}  \int_{\mathbb{R}^d} \mathd x_i
\end{equation}
and
\begin{equation}
  S (k, x) \assign \sum_{i = 0}^{N - 1} \Delta t \left[ k_i \left( \frac{x_{i
  + 1} - x_i}{\Delta t} \right) - {\color[HTML]{000000}\hat{r} (k_i, x_i,
  t_i)} \right] . \label{eq:action}
\end{equation}
If we recognize $(x_{i + 1} - x_i) / \Delta t$ as the velocity $\dot{x}
(t_i)$, then $S (k, x)$ can be seen as the Legendre transform $\int [p (t)
\dot{x} (t) - H (p (t), x (t))] \mathd t$, in which
{\color[HTML]{008000}{\tmem{$k$ is analogy to momentum $p$ and
${\color[HTML]{000000}\hat{r}  (k, x, t)}$ plays the role of Hamiltonian
${\color[HTML]{000000}H (p, x, t)}$}}}.

\section{From Integral to Differential}\label{section:Expanding Kernel as
Generalized Function}

Given a ``test function'' $\psi$ (TODO: which space does $\psi$ locates in?),
we can calculate its ``inner product'' with time evolution (equation
\ref{equ:superposition}), as
\[ \int_{\mathbb{R}^d} \mathd x \psi (x)  \left[ \mathi \frac{\partial
   \varphi}{\partial t} (x, t) \right] = \int_{\mathbb{R}^d} \mathd x
   \int_{\mathbb{R}^d} \mathd y \psi (x) {\color[HTML]{000000}r (x, y, t)}
   \varphi (y, t) . \]
In the right hand side, Taylor expanding $\psi$ at $y$ makes
\[ \int_{\mathbb{R}^d} \mathd x \psi (x)  \left[ \mathi \frac{\partial
   \varphi}{\partial t} (x, t) \right] = \sum_{n = 0}^{\infty} \frac{1}{n!} 
   \int_{\mathbb{R}^d} \mathd x \int_{\mathbb{R}^d} \mathd y
   \partial^n_{\alpha} \psi (y)  (x - y)^{\alpha} {\color[HTML]{000000}r (x,
   y, t)} \varphi (y, t), \]
where $\alpha \assign (\alpha_1, \ldots, \alpha_n)$ and $(x - y)^{\alpha}
\assign (x - y)^{\alpha_1} \times \cdots \times (x - y)^{\alpha_n}$ (recall
the abbreviations). Define the {\tmstrong{moment}}
\begin{equation}
  {\color[HTML]{000000}R_n^{\alpha}  (y, t)} \assign \int_{\mathbb{R}^d}
  \mathd x {\color[HTML]{000000}r (x, y, t)} (x - y)^{\alpha} .
  \label{eq:moment}
\end{equation}
Then, integrating over $x$ gives
\[ \int_{\mathbb{R}^d} \mathd x \psi (x)  \left[ \mathi \frac{\partial
   \varphi}{\partial t} (x, t) \right] = \sum_{n = 0}^{\infty} \frac{1}{n!} 
   \int_{\mathbb{R}^d} \mathd y \partial^n_{\alpha} \psi (y)
   {\color[HTML]{000000}R^{\alpha}_n  (y, t)} \varphi (y, t) . \]
After integration by parts and then omitting the boundary (since $\psi$ is
rapidly decreasing), we get
\[ \int_{\mathbb{R}^d} \mathd y \psi (y)  \left[ \mathi \frac{\partial
   \varphi}{\partial t} (y, t) \right] = \sum_{n = 0}^{\infty} \frac{(-
   1)^n}{n!}  \int_{\mathbb{R}^d} \mathd y \psi (y) \partial^n_{\alpha}
   [{\color[HTML]{000000}R^{\alpha}_n (y, t)} \varphi (y, t)], \]
where we have replaced $x$ by $y$ in the left hand side for making it clear.
Since $\psi$ is an arbitrary function in $\mathcal{S}_E (\mathbb{R}^d)$, we
have

\footnote{This is a quantum analogy of the Kramers--Moyal expansion in
stochastic process, which has the form
\[ \frac{\partial p}{\partial t}  (x, t) = \sum_{n = 0}^{\infty} \frac{(-
   1)^n}{n!} \partial^n_{\alpha}  [{\color[HTML]{000000}R^{\alpha}_n (x, t)} p
   (x, t)] . \]
Here, the $p$ is a probability density function that represents the stochastic
process. Remark that the only difference between the two equations is the
imaginary factor before temporal derivative.}
\begin{equation}
  \mathi \frac{\partial \varphi}{\partial t}  (x, t) = \sum_{n = 0}^{\infty}
  \frac{(- 1)^n}{n!} \partial^n_{\alpha} [{\color[HTML]{000000}R^{\alpha}_n
  (x, t)} \varphi (x, t)] . \label{eq:transexp}
\end{equation}
That is, we convert the integral equation \ref{equ:superposition} to a
differential equation. In practice, differential equation is much more
convenient than its integral correspondence for doing calculus. In summary, we
have three equivalent formalisms that describe the time evolution of
wavefunction: the original integral equation \ref{equ:superposition}, the path
integral \ref{eq:pathint}, and now the differential equation
\ref{eq:transexp}.

Interestingly, the Taylor expansion of the ``Hamiltonian''
${\color[HTML]{000000}\hat{r}  (k, y, t)}$, defined by equation
\ref{eq:hamiltonian}, also relates to the moments $R_n$s. Directly by equation
\ref{eq:hamiltonian}, we have
\[ \frac{\partial^n  {\color[HTML]{000000}\hat{r}}}{\partial k_{\alpha}}  (0,
   y, t) = \lim_{k \rightarrow 0}  \int_{\mathbb{R}^d} \mathd x \left[
   \frac{\partial^n}{\partial k_{\alpha}} \exp (- \mathi k (x - y)) \right]
   {\color[HTML]{000000}r (x, y, t)} = (- \mathi)^n  \int_{\mathbb{R}^d}
   \mathd x {\color[HTML]{000000}r (x, y, t)} (x - y)^{\alpha} . \]
The integral is recognized as ${\color[HTML]{000000}R_n^{\alpha}  (y, t)}$.
So, we find $(- \mathi)^n {\color[HTML]{000000}R_n^{\alpha}  (y, t)}$s the
Taylor coefficients of ${\color[HTML]{000000}\hat{r}  (k, y, t)}$ expanded by
$k$ at its origin. Namely,\footnote{\tmcolor{red}{But, we have to prove that
${\color[HTML]{000000}\hat{r}  (k, y, t)}$ is entire on $k$, so that its
Taylor series converges to itself.}}
\begin{equation}
  {\color[HTML]{000000}\hat{r}  (k, y, t)} = \sum_{n = 0}^{\infty} \frac{(-
  \mathi)^n}{n!} R^{\alpha}_n (y) k_{\alpha}, \label{eq:momentexpansion}
\end{equation}
where $k_{\alpha} \assign (k_{\alpha_1} \times \cdots \times k_{\alpha_n})$ as
usual. Again, the details of $S (k, x)$ (defined in equation \ref{eq:pathint})
can be completely determined by the moments $R_n$s.

So, consider the traditional Hamiltonian $\hat{r}  (x, p) = p^2 / (2 m) + V
(x)$, all $R_n$s vanish except for ${\color[HTML]{000000}R_0  (x, t)} = V (x)$
and $R_2 (x) = - 1 / m$. Equation \ref{eq:transexp} implies
\begin{equation}
  \mathi \frac{\partial \varphi}{\partial t}  (x, t) = - \frac{1}{2 m}
  \nabla^2 \varphi (x, t) + V (x) \varphi (x, t),
\end{equation}
which is exactly the Schrödinger equation (in the natural units where
Planck's constant $\hbar = 1$).

\section{Locality Truncates the Moments}\label{section:Locality Truncates the
Moments}

We then introduce the third axiom about locality, and discuss what it will
induce.

\begin{axiom}[Locality]
  \label{axiom:local}Time evolution of quantum system is local.
\end{axiom}

Shall not confuse time evolution with collapse, which is proven to be
non-local. Axiom \ref{axiom:local} claims that time evolution (equation
\ref{equ:superposition}) is local. To make this clear, we consider an example,
in which $R_n  (x, t) = c^n$ for some constant $c$, and set the dimension $d =
1$ for simplicity. Then, time evolution (equation \ref{eq:transexp}) at $x =
0$ is
\[ \mathi \frac{\partial \varphi}{\partial t}  (0, t) = \sum_{n = 0}^{\infty}
   \frac{(- c)^n}{n!}  \frac{\partial^n \varphi}{\partial x^n}  (0, t) . \]
The last expression happens to be the Taylor series of $\varphi (x, t)$ at $x
= - c$, namely $\varphi (- c, t)$. So, we conclude that $R_n  (x, t) = c^n$
for some constant $c$ implies
\[ \mathi \frac{\partial \varphi}{\partial t}  (0, t) = \varphi (- c, t) . \]
If we change the value of $\varphi$ at $x = - c$, then the time evolution at
$x = 0$ changes accordingly. It means non-locality.

In physics, a local equation generally refers to an operation on $\varphi$
which contains $\varphi$ itself and {\tmem{finite}} number of spatial
derivatives of $\varphi$, such as
\[ \mathi \frac{\partial \varphi}{\partial t}  (x, t) =\mathcal{L} \left(
   \varphi (x, t), \frac{\partial \varphi}{\partial x} (x, t),
   \frac{\partial^2 \varphi}{\partial x^2} (x, t), \ldots, \frac{\partial^n
   \varphi}{\partial x^n} (x, t) \right), \]
where $\mathcal{L}$ is an entire function. This is easy to understand because
to compute $(\partial^n \varphi / \partial x^n) (0, t)$ using numerical method
with difference $\Delta x$, only $\varphi (x, t)$ with $x \in \{0, \Delta x,
\ldots, n \Delta x\}$ are employed.\footnote{Given a smooth function $f$,
numerically evaluating $f^{(n)} (0)$ needs $f^{(n - 1)} (0)$ and $f^{(n - 1)}
(\Delta x)$. Recursively, numerically evaluating $f^{(n - 1)} (0)$ needs
$f^{(n - 2)} (0)$ and $f^{(n - 2)} (\Delta x)$, and $f^{(n - 1)} (\Delta x)$
needs $f^{(n - 2)} (\Delta x)$ and $f^{(n - 2)} (2 \Delta x)$. Repeating this
process, we find $\{f (0), \ldots, f (n \Delta x)\}$ are needed for evaluating
$f^{(n)} (0)$.} So, $(\partial \varphi / \partial t) (0, t)$ cannot
``perceive'' the $\varphi (x, t)$ outside the neighborhood $\{x : |x|
\leqslant n \Delta x\}$. Since $\Delta x$ can be arbitrarily small (but not
vanishing), the equation is local.

We have shown that, with a cut-off $N_{\tmop{cut}}$ on $R_n$s such that $R_n =
0$ for any $n > N_{\tmop{cut}}$, time evolution of wavefunction is local. And
without such a cut-off, we can construct a sequence of $R_n$s such that time
evolution is not local. Now, we are to prove that, generally, without a
cut-off, any sequence of $R_n$s (such that for any $N > 0$, there are infinite
many $R_n$s that are not vanishing), time evolution of wavefunction is
non-local. This then imports a cut-off on moments.

Before proving, we first declare a property of entire function. Consider the
Taylor expansion of $\varphi \in \mathcal{S}_E (\mathbb{R})$ at origin
\[ \varphi (x) = a_0 + a_1 x + a_2 x^2 + \cdots . \]
The information of $\varphi$ is completely encoded in the infinite sequence of
$a_n$s. This is the result of a theorem which claims that two entire functions
are equal if they agree in any interval (hence we can obtain the Taylor series
of the function within the interval). What if we only know a portion of the
infinite sequence of $a_n$s? For example, if we only know the $a_0$, then only
the value of $\varphi (x)$ at $x = 0$ is determined. Further, if we also know
the $a_1$, then we can give a good approximation to $\varphi (x)$ in a very
tiny neighborhood at $x = 0$, since $\varphi (x) = a_0 + a_1 x + \omicron
(x^2)$. Then, if we additionally know the $a_2$, then the approximation
becomes as good as before in a larger neighborhood at $x = 0$, since $\varphi
(x) = a_0 + a_1 x + a_2 x^2 + \omicron (x^3)$ and the size of neighborhood can
increase a little without growing the scale of residue. This analysis
indicates that the more $a_n$s we know, in a larger neighborhood at $x = 0$
can we properly approximate $\varphi (x)$. Our vision becomes border and
border if we know more and more $a_n$s. Until knowing the whole sequence of
$a_n$s, we realize the complete picture of $\varphi (x)$ (based on the
previous theorem about entire function).

It also indicates that, for a sufficient large $N$, we can keep $\varphi (x)$
(approximately) invariant when we tune the $a_n$s with $n > N$ while keeping
the other $a_n$s unchanged. On the other hand, based on the equation (plugging
$\varphi (x) = \sum_n a_n x^n$ into equation \ref{eq:transexp}, and collecting
all $a_m$ terms with $m < n$ into $[\cdots]$),
\[ \mathi \frac{\partial \varphi}{\partial t}  (0, t) = \sum_{n = 0}^{\infty}
   \{[\cdots] + (- 1)^n {\color[HTML]{000000}R_n (0, t)} a_n \}, \]
if there is not a cut-off to the infinite sequence of
${\color[HTML]{000000}R_n  (0, t)}$s, we can modify the $(\partial \varphi /
\partial t) (0, t)$ by tuning an $a_n$ where $n$ can be arbitrarily large.
This, however, will leave the $\varphi (x)$ around $x = 0$ invariant. It means
that the value of $\varphi$ with $x$ far away from $x = 0$ can affect the time
evolution of $\varphi$ at $x = 0$. That is, time evolution of wavefunction is
non-local. Hence, there must be a cut-off on the sequence of
${\color[HTML]{000000}R_n  (0, t)}$s. This discussion also holds for any value
of $x$ other than $x = 0$. We conclude that {\color[HTML]{008000}{\tmem{time
evolution is local if and only if there is a positive integer $N_{\tmop{cut}}$
such that $R_n = 0$ for any $n > N_{\tmop{cut}}$}}}.

\section{Hermitianity on the Moments}\label{section:Hermitianity on Moments}

Now we study the relation between $R_n$s and their complex conjugations.
Direct calculation is found tedious. Instead, we start at evaluating
$\hat{r}^{\ast}$. By conjugating the definition of $\hat{r}$ (equation
\ref{eq:hamiltonian}) and applying the Hermitianity of $r$, we get
\[ {\color[HTML]{000000}\hat{r}^{\ast}  (k, y, t)} = \int_{\mathbb{R}^d}
   \mathd x \exp (\mathi k (x - y)) {\color[HTML]{000000}r (y, x, t)} . \]
Then, applying equation \ref{eq:r-fourier} to ${\color[HTML]{000000}r (y, x,
t)}$ gives
\[ {\color[HTML]{000000}\hat{r}^{\ast}  (k, y, t)} = \int_{\mathbb{R}^d}
   \mathd k'  \int_{\mathbb{R}^d} \frac{\mathd x}{(2 \mathpi)^d} \exp \{\mathi
   (k - k') (x - y)\} {\color[HTML]{000000}\hat{r}  (k, x, t)} . \]
Taylor expanding ${\color[HTML]{000000}\hat{r}  (k', x, t)}$ by $x$ at $y$,
\[ {\color[HTML]{000000}\hat{r}^{\ast}  (k, y, t)} = \sum_{n = 0}^{\infty}
   \frac{1}{n!}  \int_{\mathbb{R}^d} \mathd k'  \frac{\partial^n 
   {\color[HTML]{000000}\hat{r}}}{\partial y^{\alpha}}  (k', y, t) \times
   \int_{\mathbb{R}^d} \frac{\mathd x}{(2 \mathpi)^d} \exp \{\mathi (k - k')
   (x - y)\}  (x - y)^{\alpha} . \]
Notice the relation
\[ \partial_{\alpha}^n \delta (\omega) = \frac{\partial^n}{\partial
   \omega^{\alpha}}  \int_{\mathbb{R}^d} \frac{\mathd z}{(2 \mathpi)^d} \exp
   (\mathi \omega z) = \mathi^n  \int_{\mathbb{R}^d} \frac{\mathd z}{(2
   \mathpi)^d} \exp (\mathi \omega z) z^{\alpha} . \]
Replacing $\omega$ by $(k - k')$ and $z$ by $(x - y)$, the last integral in
${\color[HTML]{000000}\hat{r}^{\ast}  (k, y, t)}$ becomes $(- \mathi)^n
\partial^n_{\alpha} \delta (k - k')$. Thus,
\[ {\color[HTML]{000000}\hat{r}^{\ast}  (k, y, t)} = \sum_{n = 0}^{\infty}
   \frac{(- \mathi)^n}{n!}  \int_{\mathbb{R}^d} \mathd k'  \frac{\partial^n 
   {\color[HTML]{000000}\hat{r}}}{\partial y^{\alpha}}  (k', y, t)
   \partial^n_{\alpha} \delta (k - k') . \]
By integration by parts (recalling that $\partial^n \delta$ is odd when $n$ is
odd, otherwise even) and then integrating over $k'$, we arrive at
\[ {\color[HTML]{000000}\hat{r}^{\ast}  (k, y, t)} = \sum_{n = 0}^{\infty}
   \frac{(- \mathi)^n}{n!}  \frac{\partial^{2 n} 
   {\color[HTML]{000000}\hat{r}}}{\partial k_{\alpha} \partial y^{\alpha}} 
   (k, y, t) . \]
Since $(- \mathi)^n {\color[HTML]{000000}R_n  (y, t)}$s are the coefficients
of Taylor expansion of ${\color[HTML]{000000}\hat{r}  (k, y, t)}$ by $k$
(equation \ref{eq:momentexpansion}), we find
\[ {\color[HTML]{000000}(R^{\alpha}_m)^{\ast} (y, t)} = (- \mathi)^m
   \frac{\partial^m {\color[HTML]{000000}\hat{r}^{\ast}}}{\partial k_{\alpha}}
   (0, y, t) = \sum_{n = 0}^{\infty} \frac{(- \mathi)^{m + n}}{n!} 
   \frac{\partial^n}{\partial y^{\beta}}  \frac{\partial^{m + n} 
   {\color[HTML]{000000}\hat{r}}}{\partial k_{\alpha} \partial k_{\beta}}  (0,
   y, t) . \]
Again, the $(m + n)$-th coefficients of the Taylor expansion is $(- \mathi)^{m
+ n} {\color[HTML]{000000}R_{m + n}  (y, t)}$, so
\[ {\color[HTML]{000000}(R^{\alpha}_m)^{\ast} (y, t)} = \sum_{n = 0}^{\infty}
   \frac{(- \mathi)^{m + n}}{n!}  \frac{\partial^n}{\partial x^{\beta}}  [(-
   \mathi)^{m + n} R^{\alpha \beta}_{m + n} (y)] = \sum_{n = 0}^{\infty}
   \frac{(- 1)^{m + n}}{n!} \partial^n_{\beta} {\color[HTML]{000000}R^{\alpha
   \beta}_{m + n}  (x, t)} . \]
Recall that $R_n = 0$ for any $n > N_{\tmop{cut}}$ (section
\ref{section:Locality Truncates the Moments}), we finally arrive at
\begin{equation}
  (R^{\alpha}_m)^{\ast} (x, t) = \sum_{n = 0}^{N_{\tmop{cut}} - m} \frac{(-
  1)^{m + n}}{n!} \partial^n_{\beta} {\color[HTML]{000000}R^{\alpha \beta}_{m
  + n}  (x, t)} .
\end{equation}
It relates the moments $R_n$s to their complex conjugations.

For example, when $N_{\tmop{cut}} = 2$, we have
${\color[HTML]{000000}R_0^{\ast}  (x, t)} = {\color[HTML]{000000}R_0  (x, t)}
- \partial_{\alpha} {\color[HTML]{000000}R^{\alpha}_1  (x, t)} + (1 / 2)
\partial^2_{\alpha} R^{\alpha}_2  (x, t)$, $(R_1^{\alpha})^{\ast} (x, t) = -
{\color[HTML]{000000}R_1^{\alpha}  (x, t)} + \partial_{\beta} R^{\alpha
\beta}_2  (x, t)$, and $(R_2^{\alpha})^{\ast} (x, t) =
{\color[HTML]{000000}R^{\alpha}_2  (x, t)}$, that is, $R_2$ is real.
Specially, we have
\[ {\color[HTML]{000000}R_{N_{\tmop{cut}}}^{\ast}  (x, t)} = (-
   1)^{N_{\tmop{cut}}} R_{N_{\tmop{cut}}} (x) . \]
That is, $R_{N_{\tmop{cut}}}$ is real when $N_{\tmop{cut}}$ is even, otherwise
purely imaginary.

\section{Galilean Invariance Selects $N_{\tmop{cut}} =
2$}\label{section:Galilean Invariance Selects N=2}

Hermitianity (section \ref{section:Hermitianity on Moments}) forces
$N_{\tmop{cut}}$ to be even. Locality (section \ref{section:Locality Truncates
the Moments}) guarantees a finite $N_{\tmop{cut}}$. The physical value
$N_{\tmop{cut}} = 2$ is not fixed by these general principles alone; it must
be deduced from an external physical constraint. We now show that Galilean
invariance -- the requirement that the laws of physics take the same form in
all inertial frames -- forces $N_{\tmop{cut}} \le 2$ and fixes $R_2 = - 1 /
m$.

\subsection{Galilean Transformation of the Wavefunction}

Consider an inertial frame $(x, t)$ and another frame
$(x', t')$ moving at velocity $v$ relative to $(x,t)$, so that
\begin{equation}
  x' = x - vt, \qquad t' = t. \label{eq:galilean-coords}
\end{equation}
In quantum mechanics, the wavefunction is not a scalar: it transforms with a
phase factor (see e.g. Ballentine, {\tmem{Quantum Mechanics}}, Ch. 4).
\begin{equation}
  \varphi' (x', t') = e^{if (x, t)} \varphi (x, t),
  \label{eq:galilean-transform}
\end{equation}
where $f (x, t)$ is a real-valued phase to be determined.

\subsection{Invariance of the Time Evolution Equation}

In the frame $(x,t)$, the time evolution equation is (\ref{eq:transexp})
\[ i \frac{\partial \varphi}{\partial t}  (x, t) = H (p, x ; t) \varphi (x,
   t), \qquad H (p, x ; t) \assign \sum_{n = 0}^{N_{\tmop{cut}}} \frac{(-
   i)^n}{n!}
   p_{\alpha}^n {\color[HTML]{000000}R^{\alpha}_n  (x, t)} , \]
where operator $p \assign - \mathi \partial_x$. In the frame $(x',t')$, the same
equation must hold:
\[ i \frac{\partial \varphi'}{\partial t'}  (x', t') = H (p', x' ; t')
   \varphi' (x', t'), \]
with the same functional form of $H$. Using the transformation
(\ref{eq:galilean-coords}) and (\ref{eq:galilean-transform}), the left-hand
side becomes

\begin{align}
  i \frac{\partial \varphi'}{\partial t'}  (x', t') & = i
  \frac{\partial}{\partial t}  [e^{if (x, t)} \varphi (x - vt, t)] \nonumber\\
  & = e^{if}  [- \partial_t f (x, t) + vp + H (p ; x, t)] \varphi (x - vt, t)
  . 
\end{align}

The right-hand side, after translating the momentum operator $p' = - i
\partial_{x'} = p$, acts on $\varphi'$ as
\[ H (p', x' ; t') \varphi' (x', t') = e^{if (x, t)} H (p + \partial_x f (x,
   t) ; x, t) \varphi (x - vt, t) . \]
Cancelling the common factor $e^{if}$ and recalling that $\varphi (x - vt, t)$
is arbitrary, we obtain the operator equality
\begin{equation}
  - \partial_t f (x, t) + v^{\alpha} p_{\alpha} + H (p ; x, t) = H (p + \nabla
  f (x, t) ; x, t) . \label{eq:galilean-operator}
\end{equation}
Let $g_{\alpha}  (x, t) \assign \partial_{\alpha} f (x, t)$. The similarity
transformation (Lemma \ref{eq:sim-transform}) shifts the momentum operator:
\[ e^{- \mathi F} p_{\alpha} e^{\mathi F} = p_{\alpha} + g_{\alpha}, \qquad
   \partial_{\alpha} F = g_{\alpha} . \]
   TODO: why not using $f$ instead of $F$? After all, they are the same.
By the corollary (\ref{eq:cor-identity}), $(p + g)_{\alpha}^n = e^{- \mathi F}
p_{\alpha}^n e^{\mathi F}$. Expanding in powers of $p$,
\[ (p + g)_{\alpha_1} \cdots (p + g)_{\alpha_n} = p_{\alpha_1} \cdots
   p_{\alpha_n} + ng_{(\alpha_1} p_{\alpha_2} \cdots p_{\alpha_n)}
   +\mathcal{O} (p^{n - 2}), \]
where $(\cdummy)$ denotes symmetrisation. Substituting into $H (p + g)$,

\begin{align}
  H (p + g ; x, t) & = \sum_{n = 0}^{N_{\tmop{cut}}} \frac{(- \mathi)^n}{n!}
  (p + g)_{\alpha}^n {\color[HTML]{000000}R^{\alpha}_n  (x, t)} \nonumber\\
  & = \sum_{n = 0}^{N_{\tmop{cut}}} \frac{(- \mathi)^n}{n!}
  [p_{\alpha}^n + ng_{\alpha_1}
  p_{\alpha_2} \cdots p_{\alpha_n} +\mathcal{O}(p^{n - 2})] {\color[HTML]{000000}R^{\alpha}_n  (x, t)} . 
  \label{eq:Hpg-expand}
\end{align}

TODO: I have fixed the typo in $H(p,x,t)$, where $R_n^{\alpha}$ is the the right of $p_{\alpha}$. Then continue the fixing.

\paragraph{Leading operator orders force $N_{\tmop{cut}} \le 2$.} First
consider $N_{\tmop{cut}} \ge 3$. The $v^{\alpha} p_{\alpha}$ term on the
left-hand side of (\ref{eq:galilean-operator}) contributes at order $p^1$,
which is strictly lower than $p^{N_{\tmop{cut}} - 1}$ (since $N_{\tmop{cut}} -
1 \ge 2$). Hence for the matching of $p^{N_{\tmop{cut}} - 1}$ coefficients,
the $v^{\alpha} p_{\alpha}$ term is irrelevant. The left-hand side gives
\[ \frac{(- \mathi)^{N_{\tmop{cut}} - 1}}{(N_{\tmop{cut}} - 1) !}
   R_{N_{\tmop{cut}} - 1}^{\alpha_2 \cdots \alpha_{N_{\tmop{cut}}}} . \]
On the right-hand side, the $n = N_{\tmop{cut}}$ term of (\ref{eq:Hpg-expand})
contributes an extra piece:

\begin{align*}
  & \frac{(- \mathi)^{N_{\tmop{cut}}}}{N_{\tmop{cut}} !}
  R_{N_{\tmop{cut}}}^{\alpha_1 \cdots \alpha_{N_{\tmop{cut}}}} \cdot
  N_{\tmop{cut}}  \hspace{0.17em} g_{\alpha_1} p_{\alpha_2} \cdots
  p_{\alpha_{N_{\tmop{cut}}}}\\
  & \qquad + \frac{(- \mathi)^{N_{\tmop{cut}} - 1}}{(N_{\tmop{cut}} - 1) !}
  R_{N_{\tmop{cut}} - 1}^{\alpha_2 \cdots \alpha_{N_{\tmop{cut}}}} .
\end{align*}

Equating the two sides and cancelling the common $R_{N_{\tmop{cut}} - 1}$ term
yields
\begin{equation}
  \frac{(- \mathi)^{N_{\tmop{cut}}}}{(N_{\tmop{cut}} - 1) !}
  R_{N_{\tmop{cut}}}^{\alpha_1 \cdots \alpha_{N_{\tmop{cut}}}} 
  \hspace{0.17em} g_{\alpha_1} = 0. \label{eq:constraint-N}
\end{equation}
A non-trivial boost requires $g = \nabla f \neq 0$, therefore
$R_{N_{\tmop{cut}}} = 0$. By descent, the highest non-zero moment cannot
exceed $N_{\tmop{cut}} - 1$. Repeating the argument shows that $R_n = 0$ for
all $n \ge 3$. Hence
\[ \boxed{N_{\tmop{cut}} \le 2} . \]
If $N_{\tmop{cut}} = 0$, equation (\ref{eq:galilean-operator}) reduces to $-
\partial_t f + v^{\alpha} p_{\alpha} + R_0 = R_0$, which forces $v = 0$ and
contradicts a non-trivial boost. Hermitianity (Section
\ref{section:Hermitianity on Moments}) forces $N_{\tmop{cut}}$ to be even,
leaving $\boxed{N_{\tmop{cut}} = 2}$ as the only admissible case.

\paragraph{Analysis of $N_{\tmop{cut}} = 2$.} With $N_{\tmop{cut}} = 2$ the
Hamiltonian reads
\[ H (p ; x, t) = R_0  (x, t) - \mathi R_1^{\alpha}  (x, t) p_{\alpha} -
   \frac{1}{2} R_2^{\alpha \beta}  (x, t) p_{\alpha} p_{\beta} . \]
Insert this into (\ref{eq:galilean-operator}) and expand the shifted operators
using the explicit formulas from Appendix \ref{appendix:operator-identity}:

\begin{align*}
  (p + g)_{\alpha} & = p_{\alpha} + g_{\alpha},\\
  (p + g)_{\alpha}  (p + g)_{\beta} & = p_{\alpha} p_{\beta} + 2 g_{(\alpha}
  p_{\beta)} + g_{\alpha} g_{\beta} - \mathi \partial_{\alpha} g_{\beta},
\end{align*}

where $g_{(\alpha} p_{\beta)} \assign \tfrac{1}{2} (g_{\alpha} p_{\beta} +
g_{\beta} p_{\alpha})$. Since $R_2^{\alpha \beta}$ is symmetric, $R_2^{\alpha
\beta} g_{(\alpha} p_{\beta)} = R_2^{\alpha \beta} g_{\alpha} p_{\beta}$.

Matching the coefficients of $p_{\beta}$ gives
\begin{equation}
  v^{\beta} - \mathi R_1^{\beta} = - \mathi R_1^{\beta} - R_2^{\alpha \beta}
  g_{\alpha} \quad \Longrightarrow \quad v^{\beta} = - R_2^{\alpha \beta}
  g_{\alpha} . \label{eq:gal-p1}
\end{equation}
For a spinless particle in the absence of a magnetic field the vector field
$R_1$ vanishes (there is no preferred direction). TODO: explain this idea, or not using this idea (hence $R_1 \neq 0$). Setting $R_1 = 0$ and
matching the $p^0$ terms yields
\begin{equation}
  - \partial_t f = - \frac{1}{2} R_2^{\alpha \beta}  (g_{\alpha} g_{\beta} -
  \mathi \partial_{\alpha} g_{\beta}) . \label{eq:gal-p0}
\end{equation}
Because the right-hand side of (\ref{eq:gal-p1}) involves only $t$-dependent
quantities, $g = \nabla f$ depends only on $t$, so $\partial_{\alpha}
g_{\beta} = 0$. Equation (\ref{eq:gal-p0}) therefore simplifies to
\[ \partial_t f = \frac{1}{2} R_2^{\alpha \beta} g_{\alpha} g_{\beta} . \]
Integrating, and using $\partial_{\alpha} f = g_{\alpha}$, we obtain
\begin{equation}
  f (x, t) = g_{\alpha} x^{\alpha} + \frac{1}{2} R_2^{\alpha \beta} g_{\alpha}
  g_{\beta}  \hspace{0.17em} t + \text{const} . \label{eq:f-general}
\end{equation}
\paragraph{Identification of $R_2$ and the phase factor.} Solving
(\ref{eq:gal-p1}) for $g$ requires inverting $R_2^{\alpha \beta}$. In the
isotropic case $R_2^{\alpha \beta} = r \hspace{0.17em} \delta^{\alpha \beta}$
we obtain $g_{\alpha} = - v_{\alpha} / r$. Substituting into
(\ref{eq:f-general}),
\[ f (x, t) = - \frac{v_{\alpha} x^{\alpha}}{r} + \frac{v^2}{2 r} t. \]
Comparison with the standard transformation law (Ballentine, {\tmem{Quantum
Mechanics}}, Ch.~4) fixes $r = - 1 / m$, i.e.
\[ \boxed{R_2^{\alpha \beta} = - \frac{1}{m} \delta^{\alpha \beta}}, \qquad
   \boxed{f (x, t) = mv_{\alpha} x^{\alpha} - \frac{1}{2} mv^2 t} . \]
Thus Galilean invariance forces $N_{\tmop{cut}} = 2$, $R_2 = - 1 / m$, and
$R_1 = 0$ (in the absence of a magnetic field). The remaining moment $R_0  (x,
t)$ is identified with the potential energy $V (x, t)$, and equation
(\ref{eq:transexp}) reduces to the standard Schrödinger equation.

\section{From Quantum to Classical}\label{section:From Quantum to Classical}

Consider the length scale of quantum world, denoted by $\lambda_Q$, and that
of classical world, denoted by $\lambda_C$. We must have $\lambda_Q \ll
\lambda_C$. In the units where $L = T$, consider the re-scaling $x' \assign x
/ \lambda_Q$ and $t' \assign t / \lambda_Q$, and use $(x', t')$-coordinates in
the three equivalent formalisms of time evolution (equations
\ref{equ:superposition}, \ref{eq:pathint}, and \ref{eq:transexp}). By
converting back to $(x, t)$-coordinates, for example, equation
\ref{eq:transexp} turns to be
\begin{equation}
  \mathi \lambda_Q  \frac{\partial \varphi}{\partial t}  (x, t) = \sum_{n =
  0}^{\infty} \frac{(- \lambda_Q)^n}{n!} \partial^n_{\alpha}
  [{\color[HTML]{000000}R^{\alpha}_n (x, t)} \varphi (x, t)],
  \label{eq:transexp-hbar}
\end{equation}
where the $R_n$s depend on $\lambda_Q$ in such a way that
${\color[HTML]{000000}R_n  (x, t)} = {\color[HTML]{000000}R'_n  (x /
\lambda_Q, t / \lambda_Q)}$ for some $\lambda_Q$-independent $R'_n$. Also for
path integral, we have
\begin{equation}
  \varphi (x_N, t + N \Delta t) = \int D (k, x) \exp \left(
  \frac{\mathi}{\lambda_Q} S (k, x) \right) \varphi (x_0, t) + \omicron
  (\Delta t), \label{eq:pathint-hbar}
\end{equation}
where, by inserting the expansion \ref{eq:momentexpansion} into the $S (k, x)$
given by equation \ref{eq:action},
\[ S (k, x) = \sum_{i = 0}^{N - 1} \Delta t \left[ k_i \left( \frac{x_{i + 1}
   - x_i}{\Delta t} \right) - \sum_{n = 0}^{N_{\tmop{cut}}} \frac{(-
   \mathi)^n}{n!} {\color[HTML]{000000}R^{\alpha}_n (x_i, t_i)} (k_i)_{\alpha}
   \right] . \]
The $1 / \lambda_Q$ in equation \ref{eq:pathint-hbar} is extracted from the
$\Delta t$ factor in $S (k, x)$. Hence, $\lambda_Q$ is absent in $S (k, x)$
except for in the $R_n$s.

The quantum scales for both space and time have the order of $\lambda_Q$. And
we cannot observe the quantum world directly using naked eyes because
$\lambda_Q$ is an extremely small number. Suppose that $R_n  (x, t)$ are
macroscopic, indicating that $\| \partial_x R_n \| \sim \| \partial_t R_n \|
\sim 1 / \lambda_C$. But since $R_n  (x, t) = R'_n  (x / \lambda_Q, t /
\lambda_Q)$, it implies that $\| \partial_x R_n' \| \sim \| \partial_t R_n' \|
\sim \lambda_Q / \lambda_C$. In other words, $R'_n$s are slowly varying in
quantum scale. Stationary phase approximation (see appendix
\ref{appendix:Stationary Phase Approximation}) claims that, as $\lambda_Q \ll
\lambda_C$, path integral is dominated by the ``path'' that has ``stationary
phase'', that is the $(k_{\star}, x_{\star})$ such that
\[ \frac{\partial S}{\partial (k_i)_{\alpha}}  (k_{\star}, x_{\star}) =
   \frac{\partial S}{\partial x_i^{\alpha}}  (k_{\star}, x_{\star}) = 0, \]
for each $i \in \{0, \ldots, N - 1\}$ and $\alpha \in \{1, \ldots, d\}$ (see
appendix \ref{appendix:Stationary Phase Approximation}). Conversely, this
system of equations, named {\tmstrong{equations of motion}}, furnishes the
classical (discrete) trajectory of particles, $(k_{\star}, x_{\star})$. In the
continuous version, the discrete index $i$ is replaced by a continous variable
$t \in [0, T]$ for some $T \in \mathbb{R}$. Then, $k_i$ and $x_i$ come to be
functions $k (t)$ and $x (t)$ respectively. Accordingly, $S (k, x)$ becomes a
functional, as
\[ S [k, x] = \int_0^T \mathd t \left[ k (t) \dot{x} (t) - \sum_{n =
   0}^{N_{\tmop{cut}}} \frac{(- \mathi)^n}{n!} R^{\alpha}_n (x (t), t)
   k_{\alpha} (t) \right], \]
where we use square brackets instead of parentheses for emphasizing that $S$
is a functional. The equations of motion become
\[ \frac{\delta S}{\delta k_{\alpha} (t)}  [k_{\star}, x_{\star}] =
   \frac{\delta S}{\delta x^{\alpha} (t)}  [k_{\star}, x_{\star}] = 0, \]
for each $t$ and index $\alpha$. It is recognized as the Hamiltonian equations
in classical mechanics. Indeed, if denote
\[ H (k, x) \assign \sum_{n = 0}^{N_{\tmop{cut}}} \frac{(- \mathi)^n}{n!}
   {\color[HTML]{000000}R^{\alpha}_n  (x, t)} k_{\alpha}, \]
then by variation, we have
\[ \frac{\delta S}{\delta k (t)} = \dot{x} (t) - \frac{\partial H}{\partial k}
   (k (t), x (t)) = 0 \Rightarrow \dot{x} (t) = \frac{\partial H}{\partial k}
   (k (t), x (t)), \]
and
\[ \frac{\delta S}{\delta x (t)} = - \dot{k} (t) - \frac{\partial H}{\partial
   x}  (k (t), x (t)) = 0 \Rightarrow \dot{k} (t) = - \frac{\partial
   H}{\partial x}  (k (t), x (t)) . \]
The right hand sides of $\Rightarrow$ are exactly the Hamiltonian equations.

\appendix\chapter{Taylor Reminder}

Given a function $f \in \mathcal{C}^n (\mathbb{R}^d)$, we are to calculate the
reminder $h_n (x, x_0)$ in the Taylor expansion at $x_0$
\[ f (x) = f (x_0) + \sum_{k = 1}^n \frac{1}{k!}  (\partial_{\alpha_1} \cdots
   \partial_{\alpha_k} f) (x_0)  \prod_{i = 1}^k (x^{\alpha_i} -
   x_0^{\alpha_i}) + h_{n + 1} (x, x_0) . \]
We begin at a variation of the fundamental theorem of calculus
\[ f (x) - f (x_0) = \int_0^1 \mathd t (\partial_{\alpha_1} f)  (x + t (x_0 -
   x))  (x^{\alpha_1} - x_0^{\alpha_1}) . \]
Integration by parts gives
\[ f (x) - f (x_0) = (\partial_{\alpha_1} f) (x_0)  (x^{\alpha_1} -
   x_0^{\alpha_1}) + \frac{1}{2}  \int_0^1 \mathd (t^2)  (\partial_{\alpha_1}
   \partial_{\alpha_2} f)  (x + t (x_0 - x))  (x^{\alpha_1} - x_0^{\alpha_1}) 
   (x^{\alpha_2} - x_0^{\alpha_2}) . \]
Applying integration by parts on the integral again, we get

\begin{align}
  f (x) - f (x_0) = & (\partial_{\alpha_1} f) (x_0)  (x^{\alpha_1} -
  x_0^{\alpha_1}) \nonumber\\
  + & \frac{1}{2!}  (\partial_{\alpha_1} \partial_{\alpha_2} f) (x_0) 
  (x^{\alpha_1} - x_0^{\alpha_1})  (x^{\alpha_2} - x_0^{\alpha_2}) \nonumber\\
  + & \frac{1}{3!}  \int_0^1 \mathd (t^3)  (\partial_{\alpha_1}
  \partial_{\alpha_2} \partial_{\alpha_3} f)  (x + t (x_0 - x))  (x^{\alpha_1}
  - x_0^{\alpha_1})  (x^{\alpha_2} - x_0^{\alpha_2})  (x^{\alpha_3} -
  x_0^{\alpha_3}) . \nonumber
\end{align}

Repeating this process gives

\begin{align}
  f (x) - f (x_0) = & (\partial_{\alpha_1} f) (x_0)  (x^{\alpha_1} -
  x_0^{\alpha_1}) \nonumber\\
  + & \frac{1}{2!}  (\partial_{\alpha_1} \partial_{\alpha_2} f)  (x^{\alpha_1}
  - x_0^{\alpha_1})  (x^{\alpha_2} - x_0^{\alpha_2}) \nonumber\\
  + & \cdots \nonumber\\
  + & \frac{1}{n!}  \int_0^1 \mathd (t^n)  (\partial_{\alpha_1} \cdots
  \partial_{\alpha_n} f)  (x + t (x_0 - x))  (x^{\alpha_1} - x_0^{\alpha_1})
  \cdots (x^{\alpha_n} - x_0^{\alpha_n}) . \nonumber
\end{align}

So, the Taylor reminder reads
\begin{equation}
  h_n (x, x_0) = \frac{1}{n!}  \prod_{i = 1}^n (x^{\alpha_i} - x_0^{\alpha_i})
  \times \int_0^1 \mathd (t^n)  (\partial_{\alpha_1} \cdots
  \partial_{\alpha_n} f)  (x + t (x_0 - x)) .
\end{equation}
\chapter{Rapidly Decreasing Functions are Dense in Square-Integrable Space
(TODO)}\label{appendix:dense}

A function $f : \mathbb{R}^d \rightarrow \mathbb{C}$ is square-integrable if
\[ \int_{\mathbb{R}^d} \mathd x |f (x) |^2 \assign \int_{\mathbb{R}^d} \mathd
   xf^{\ast} (x) f (x) < \infty . \]
This is a improper integral, thus defined by
\[ \int_{\mathbb{R}^d} \mathd x |f (x) |^2 = \lim_{R \rightarrow \infty} 
   \int_{B (R)} \mathd x |f (x) |^2, \]
where $B (R) \assign [- R, R]^d$ denotes the $d$-dimensional ``box''. In other
words, for any $\varepsilon > 0$, there exists $R_{\star} > 0$ and $\delta >
0$, such that for any $R > R_{\star}$ and $0 < \Delta x < \delta$,
\[ \left| \int_{\mathbb{R}^d} \mathd x|f (x) |^2 - \int_{B (R)} \mathd x|f (x)
   |^2 \right| < \varepsilon . \]
So, for any $\varepsilon > 0$, we can construct a compact supported function
$g_R : \mathbb{R}^d \rightarrow \mathbb{C}$, which consists with $f$ within
the box $B (R)$ and vanishes outside.

\

\hrulefill

\

Consider the convolution
\[ f_n (x) = (\delta_n \ast f) (x) \assign \int_{\mathbb{R}^d} \mathd y
   \delta_n  (x - y) f (y), \]
where $n$ is a positive integer and
\[ \delta_n (x) \assign \frac{1}{\left( \sqrt{2 \mathpi} n \right)^d} \exp
   \left( - \frac{x^2}{2 n} \right), \]
that is, the probabilistic density function of normal distribution with zero
mean and variance $n$. By young's inequality for convolution, we have
\[ \| \delta_n \ast f\|_{\infty} \leqslant \| \delta_n \|_2 \|f\|_2 . \]

\chapter{Action Is Real in the Classical
Limit}\label{appendix:real-action-cl}

The stationary phase approximation (Appendix
\ref{appendix:Stationary Phase Approximation}) requires a real-valued
$S$. We show that $\Im S = O (\lambda_Q)$, so in the classical limit
$\lambda_Q \to 0$ the action is strictly real.

\paragraph{Hermitianity and the imaginary part of $\hat{r}$.}
Section \ref{section:Hermitianity on Moments} derived, from the
Hermitianity $r^{\ast} (x, y, t) = r (y, x, t)$,
\begin{equation}
  \hat{r}^{\ast} (k, y, t) = \sum_{n = 0}^{N_{\tmop{cut}}}
  \frac{(- \mathi)^n}{n!} \frac{\partial^{2 n} \hat{r}}{\partial k_{\alpha}
  \partial y^{\alpha}} (k, y, t) . \label{eq:rstar-sum}
\end{equation}
Separating the $n = 0$ term ($= \hat{r}$) and using $2 \mathi \Im \hat{r}
= \hat{r} - \hat{r}^{\ast}$, we obtain
\begin{equation}
  \Im \hat{r} (k, y, t) = \frac{\mathi}{2} \sum_{n = 1}^{N_{\tmop{cut}}}
  \frac{(- \mathi)^n}{n!} \frac{\partial^{2 n} \hat{r}}{\partial k_{\alpha}
  \partial y^{\alpha}} (k, y, t) . \label{eq:Imr-sum}
\end{equation}
Every term in this sum contains at least one $y$-derivative.

\paragraph{Slow-variation estimate.}
Section \ref{section:From Quantum to Classical} introduces the scaling
$R_n (x, t) = R'_n (x / \lambda_Q, t / \lambda_Q)$ where $R'_n$ is
independent of $\lambda_Q$ and varies only on the macroscopic scale.  In
quantum units (where $\lambda_Q$ is the unit of length), the macroscopic
scale $\lambda_C$ is large, so
\[ \| \partial_y R_n \| = \frac{1}{\lambda_Q} \| \partial_{y'} R'_n \|
   \sim \frac{1}{\lambda_C} \ll 1 . \]

Since $\hat{r}$ in (\ref{eq:Imr-sum}) is a polynomial in $k$ whose
coefficients are the moments $R_n$ and their $y$-derivatives, each
$\partial_y$ acting on $\hat{r}$ brings down factors of $\partial_y R_n$.
Hence
\begin{equation}
  \Im \hat{r} (k, y, t) = O \bigl( 1 / \lambda_C \bigr) .
  \label{eq:Imr-order}
\end{equation}

\paragraph{Imaginary part of the action.}
In the discretised action (\ref{eq:action}),
\[ S (k, x) = \sum_{i = 0}^{N - 1} \Delta t \Bigl[ (k_i)_{\alpha}
   \frac{x_{i + 1}^{\alpha} - x_i^{\alpha}}{\Delta t} -
   \hat{r} (k_i, x_i, t_i) \Bigr], \]
the imaginary part is
\[ \Im S = - \sum_i \Delta t \; \Im \hat{r} (k_i, x_i, t_i)
   = O (T / \lambda_C), \qquad T \assign N \Delta t . \]

In quantum units $\lambda_C \gg 1$, so $\Im S \ll 1$.  The path integral
factor is
\[ \exp \bigl( \mathi S / \lambda_Q \bigr) = \exp \bigl( \mathi S_R
   / \lambda_Q \bigr) \; \exp \bigl( - \Im S / \lambda_Q \bigr) . \]
Since $\Im S / \lambda_Q = O (T / (\lambda_Q \lambda_C)) = O (T)$ (using
$\lambda_Q \lambda_C \sim 1$ in the physical regime), the damping factor
remains $O (1)$.  In the classical limit $\lambda_Q \to 0$ (equivalently
$\lambda_C \to \infty$), $\Im S \to 0$ and the action is strictly real.
The stationary phase approximation is therefore applicable in this
limit.

\chapter{Conditions for Preserving Schwartz
Space}\label{appendix:schwartz-preserve}

We address the question raised in Section \ref{section:Wavefunction Is
Rapidly Decreasing}: what condition guarantees that the time evolution
$i \partial_t \varphi = \int r \varphi$ keeps $\varphi (\cdot, t)$ in the
Schwartz space $\mathcal{S} (\mathbb{R}^d)$ for all $t$?

\paragraph{Differential form of the evolution.}
The superposition principle (Section \ref{section:Superposition Principle
and Time Evolution}) and locality (Section \ref{section:Locality Truncates
the Moments}) together imply that the kernel $r$ can be written as a
finite-order differential operator:
\begin{equation}
  \mathi \frac{\partial \varphi}{\partial t}  (x, t) = \sum_{n =
  0}^{N_{\tmop{cut}}} \frac{(- 1)^n}{n!} \partial_{\alpha}^n
  [{\color[HTML]{000000}R^{\alpha}_n  (x, t)} \varphi (x, t)] .
  \label{eq:diff-evolution}
\end{equation}
This is a linear partial differential equation of Schrödinger type. The
question reduces to: when does such an operator $\mathcal{L} \varphi \equiv
\sum_n \frac{(-1)^n}{n!} \partial^n [R_n \varphi]$ preserve the Schwartz
space?

\paragraph{Sufficient condition on the coefficients.}
The Schwartz space $\mathcal{S} (\mathbb{R}^d)$ is a Fréchet space
topologised by the seminorms
\[ \| f \|_{k, m} \assign \sup_{|\alpha| \le k, |\beta| \le m} \;
   \sup_{x \in \mathbb{R}^d} | x^{\alpha} \partial^{\beta} f (x) | . \]
A linear differential operator maps $\mathcal{S}$ continuously into
$\mathcal{S}$ if its coefficients are {\tmstrong{smooth and slowly
increasing}} (also called {\tmstrong{tempered}}).  Explicitly, for each
$n$ and each multi-index $\alpha$, the function $R_n^{\alpha} (x, t)$ must
be $C^{\infty}$ in $x$ and satisfy
\begin{equation}
  \sup_{x \in \mathbb{R}^d} \bigl| \partial_x^{\beta} R_n^{\alpha} (x, t)
  \bigr| \le C_{\beta} (1 + |x|)^{M_{\beta}}
  \qquad (\forall \beta)
  \label{eq:tempered-coeff}
\end{equation}
for some constants $C_{\beta}, M_{\beta}$ depending on $\beta$ but not on
$x$.  This condition is standard in the theory of partial differential
equations: it ensures that the operator, and all its compositions with
polynomials and derivatives, remain bounded on $\mathcal{S}$.

\paragraph{Preservation of Schwartz space.}
When (\ref{eq:tempered-coeff}) holds for all $n \le N_{\tmop{cut}}$, the
right-hand side of (\ref{eq:diff-evolution}) belongs to $\mathcal{S}$
whenever $\varphi (\cdot, t) \in \mathcal{S}$.  The time derivative
$\partial_t \varphi$ is then also in $\mathcal{S}$.  By the standard theory
of linear evolution equations in Fréchet spaces (or, more concretely, by
uniqueness of smooth solutions to the Schrödinger equation with a tempered
Hamiltonian), we obtain
\[ \varphi (\cdot, 0) \in \mathcal{S} (\mathbb{R}^d)
   \quad\Longrightarrow\quad
   \varphi (\cdot, t) \in \mathcal{S} (\mathbb{R}^d) \quad \forall t . \]

\paragraph{Physical examples.}
The familiar Schrödinger Hamiltonian $H = - \nabla^2 / (2 m) + V (x)$
corresponds to $N_{\tmop{cut}} = 2$ with $R_2 = - 1 / m$ (constant),
$R_1 = 0$, and $R_0 = V (x)$.  Condition (\ref{eq:tempered-coeff}) is
satisfied whenever $V$ and its derivatives grow at most polynomially.
This includes the harmonic oscillator ($V \propto x^2$), anharmonic
oscillators, power-law potentials, and many others.

\paragraph{Consistency with Fourier analysis.}
The notes perform Fourier transforms on the kernel $r (x, y, t)$ with
respect to its first argument (Section \ref{section:Path Integral
Formalism}, equations \ref{eq:r-fourier} and \ref{eq:hamiltonian}).
The Fourier theory on which these operations rely is reviewed in
Appendix \ref{appendix:fourier-on-sprime}.

\chapter{Fourier Transform on $\mathcal{S}$ and
$\mathcal{S}'$}\label{appendix:fourier-on-sprime}

Fourier transform is used throughout this work on the kernel
$r (x, y, t)$ and its Fourier image $\hat{r} (k, y, t)$.  We summarise
here the functional-analytic facts that guarantee these operations are
well-defined.

\paragraph{The Schwartz space $\mathcal{S} (\mathbb{R}^d)$.}
A function $f : \mathbb{R}^d \to \mathbb{C}$ belongs to the Schwartz space
$\mathcal{S} (\mathbb{R}^d)$ if it is $C^{\infty}$ and decays faster than
any polynomial together with all its derivatives.  The Fourier transform
\begin{equation}
  \tilde f (k) = \int_{\mathbb{R}^d} \mathd x \, e^{- \mathi k x}
  f (x), \qquad k \in \mathbb{R}^d
\end{equation}
is an {\tmstrong{automorphism}} of $\mathcal{S} (\mathbb{R}^d)$:
if $f \in \mathcal{S}$ then $\tilde f \in \mathcal{S}$, and the inverse
transform
\[ f (x) = \int_{\mathbb{R}^d} \frac{\mathd k}{(2 \mathpi)^d} \,
   e^{\mathi k x} \tilde f (k) \]
maps $\mathcal{S}$ back onto $\mathcal{S}$.

\paragraph{Tempered distributions $\mathcal{S}' (\mathbb{R}^d)$.}
Many objects of interest---polynomials, Dirac $\delta$, derivatives of
$\delta$---are not in $\mathcal{S}$ because they do not decay at infinity.
They are, however, {\tmstrong{tempered distributions}}, i.e. continuous
linear functionals on $\mathcal{S}$.  The space of all such functionals is
denoted $\mathcal{S}' (\mathbb{R}^d)$.  Examples:
\begin{align*}
  \text{polynomial } P (k) &\longmapsto \langle P, \varphi \rangle =
  \int \mathd k \, P (k) \varphi (k),\\
  \text{Dirac } \delta (x) &\longmapsto \langle \delta, \varphi \rangle =
  \varphi (0),\\
  \text{derivative } \partial^n \delta (x) &\longmapsto
  \langle \partial^n \delta, \varphi \rangle = (- 1)^n
  \partial^n \varphi (0) .
\end{align*}
The duality pairing $\langle \cdot, \cdot \rangle$ is bilinear and extends
the usual $L^2$ inner product whenever the latter is defined.

\paragraph{Fourier transform on $\mathcal{S}'$ is also an automorphism.}
The Fourier transform is extended from $\mathcal{S}$ to $\mathcal{S}'$
{\tmstrong{by duality}}.  For a tempered distribution $T \in
\mathcal{S}'$, define $\tilde T \in \mathcal{S}'$ by
\begin{equation}
  \langle \tilde T, \varphi \rangle \assign \langle T, \tilde \varphi
  \rangle, \qquad \forall \varphi \in \mathcal{S}.
  \label{eq:ft-on-Sprime}
\end{equation}
Because $\tilde \varphi \in \mathcal{S}$ (automorphism on $\mathcal{S}$),
the right-hand side is well-defined for every test function $\varphi$.
Laurent Schwartz proved that this definition makes the Fourier transform
an {\tmstrong{automorphism}} of $\mathcal{S}'$: the map
$\mathcal{F} : \mathcal{S}' \to \mathcal{S}'$ is one-to-one, onto, and its
inverse is given by the dual of the inverse Fourier transform on
$\mathcal{S}$.  In particular, for any $T \in \mathcal{S}'$,
$\mathcal{F} T$ is again a tempered distribution and the inversion formula
holds in the sense of distributions.

\paragraph{Application to the kernel $r$.}
In the local form (Section \ref{section:Locality Truncates the Moments}),
\[ r (x, y, t) = \sum_{n = 0}^{N_{\tmop{cut}}} \frac{(- 1)^n}{n!}
   R_n^{\alpha} (y, t) \, \partial_x^n \delta (x - y) . \]
For fixed $y$ and $t$, this is a finite linear combination of $\delta$ and
its derivatives --- each term is a tempered distribution in $x$.  Hence
$r (\cdot, y, t) \in \mathcal{S}' (\mathbb{R}^d_x)$, and its Fourier
transform $\hat{r} (\cdot, y, t)$ (defined by
(\ref{eq:ft-on-Sprime})) is also a tempered distribution.  Explicitly,
\begin{equation}
  \hat{r} (k, y, t) = \sum_{n = 0}^{N_{\tmop{cut}}}
  \frac{(- \mathi)^n}{n!} R_n^{\alpha} (y, t) \, k_{\alpha},
  \label{eq:rhat-polynomial}
\end{equation}
which is a polynomial in $k$ --- a tempered distribution that is smooth
and slowly increasing.  The inverse Fourier transform recovers $r$ from
$\hat{r}$ by the same dual procedure.

Thus all Fourier operations on $r$ and $\hat{r}$ in Sections
\ref{section:Path Integral Formalism} and
\ref{section:Hermitianity on Moments} are mathematically legitimate.

\chapter{Stationary Phase Approximation}\label{appendix:Stationary Phase
Approximation}

\paragraph{Motivation.} In Section \ref{section:Path Integral Formalism} the
time evolution of the wavefunction is expressed as a path integral
\[ \varphi (x_N, t + N \Delta t) = \int D (k, x) \exp \left(
   \frac{\mathi}{\lambda_Q} S (k, x) \right) \varphi (x_0, t) + \omicron
   (\Delta t), \]
where $S (k, x)$ is given by (\ref{eq:action}). By Appendix
\ref{appendix:real-action-cl}, $\Im S = O (T / \lambda_C)$, so in the
classical limit $\lambda_C \to \infty$ (equivalently $\lambda_Q \to 0$)
the exponential is a pure oscillation.
When the quantum scale $\lambda_Q$ is much smaller than the classical scale
$\lambda_C$, this exponential oscillates rapidly and the dominant
contribution comes from the neighbourhood of stationary points of $S$.
This is the content of the stationary phase approximation (also known as
the method of steepest descent).

\paragraph{The multi-dimensional stationary phase theorem.} Let $S :
\mathbb{R}^m \to \mathbb{R}$ be a smooth function and $A : \mathbb{R}^m \to
\mathbb{C}$ be a smooth amplitude that grows at most polynomially at infinity.
Suppose $z_0 \in \mathbb{R}^m$ is a {\tmstrong{non-degenerate critical point}}
of $S$, i.e.
\[ \nabla S (z_0) = 0, \qquad \det H_S (z_0) \neq 0, \]
where $H_S (z_0)$ denotes the Hessian matrix $\partial_i \partial_j S (z_0)$.
Then, as $\lambda \to 0^+$,
\[ I (\lambda) \assign \int_{\mathbb{R}^m} \mathd z \hspace{0.27em} A (z) 
   \hspace{0.17em} e^{\mathi S (z) / \lambda} \sim \left( \frac{2 \mathpi
   \mathi \lambda}{\det H_S (z_0)} \right)^{\hspace{-0.17em} m / 2} A (z_0) 
   \hspace{0.17em} e^{\mathi S (z_0) / \lambda} . \]
If $S$ has multiple non-degenerate critical points, the asymptotic expansion
is the sum over all such points.

\paragraph{{\tmem{Proof sketch.}}} Taylor expand $S$ around $z_0$ to second
order:
\[ S (z) = S (z_0) + \frac{1}{2}  (z - z_0)^{\mathsf{T}} H_S (z_0)  (z - z_0)
   +\mathcal{O} (|z - z_0 |^3) . \]
Change variables $u = (z - z_0) / \sqrt{\lambda}$, so that
\[ I (\lambda) = \lambda^{m / 2}  \int_{\mathbb{R}^m} \mathd u \hspace{0.27em}
   A (z_0 + \sqrt{\lambda} u)  \hspace{0.17em} \exp \hspace{-0.17em} \left(
   \tfrac{\mathi}{\lambda} S (z_0) + \tfrac{\mathi}{2} u^{\mathsf{T}} H_S
   (z_0) u +\mathcal{O}(\sqrt{\lambda}) \right) . \]
Taking $\lambda \to 0$, the cubic and higher terms vanish. The Gaussian
integral is
\[ \int_{\mathbb{R}^m} \mathd u \hspace{0.27em} \exp \hspace{-0.17em} \left(
   \frac{\mathi}{2} u^{\mathsf{T}} H_S (z_0) u \right) = \frac{(2 \mathpi
   \mathi)^{m / 2}}{\sqrt{\det H_S (z_0)}}, \]
where the branch of the square root is chosen by continuity in the complex
plane. Combining the factors gives the stated result.

\paragraph{Application to the discrete path integral.} Identify the
integration variables as a single vector
\[ z = (k_0, \ldots, k_{N - 1}, x_0, \ldots, x_{N - 1}) \in \mathbb{R}^{2 Nd},
\]
so that $m = 2 Nd$. The amplitude is $A (z) = \varphi (x_0, t)$, which is
slowly varying compared with the oscillatory factor (because $R'_n$s are
slowly varying on the quantum scale; see Section \ref{section:From Quantum to
Classical}). The small parameter is $\lambda = \lambda_Q$.

Applying the stationary phase theorem directly, the dominant contribution
comes from points $(k_{\star}, x_{\star})$ satisfying
\[ \frac{\partial S}{\partial (k_i)_{\alpha}}  (k_{\star}, x_{\star}) = 0,
   \qquad \frac{\partial S}{\partial x_i^{\alpha}}  (k_{\star}, x_{\star}) =
   0, \]
for every $i \in \{0, \ldots, N - 1\}$ and $\alpha \in \{1, \ldots, d\}$.

\paragraph{Explicit equations of motion.} Using the definition
(\ref{eq:action}) of $S (k, x)$,
\[ S (k, x) = \sum_{i = 0}^{N - 1} \Delta t \left[ (k_i)_{\alpha} \frac{x_{i +
   1}^{\alpha} - x_i^{\alpha}}{\Delta t} - {\color[HTML]{000000}\hat{r} (k_i,
   x_i, t_i)} \right], \]
the stationary conditions become

\begin{align*}
  \frac{\partial S}{\partial (k_i)_{\alpha}} & = \Delta t \left( \frac{x_{i +
  1}^{\alpha} - x_i^{\alpha}}{\Delta t} \right) - \Delta t \hspace{0.17em}
  \frac{\partial \hat{r}}{\partial (k_i)_{\alpha}}  (k_i, x_i, t_i) = 0,\\
  \frac{\partial S}{\partial x_i^{\alpha}} & = \Delta t [(k_i)_{\alpha} -
  (k_{i - 1})_{\alpha}] - \Delta t \hspace{0.17em} \frac{\partial
  \hat{r}}{\partial x_i^{\alpha}}  (k_i, x_i, t_i) = 0.
\end{align*}

Rearranging,
\[ \frac{x_{i + 1}^{\alpha} - x_i^{\alpha}}{\Delta t} = \frac{\partial
   \hat{r}}{\partial (k_i)_{\alpha}}, \qquad \frac{(k_i)_{\alpha} - (k_{i -
   1})_{\alpha}}{\Delta t} = - \frac{\partial \hat{r}}{\partial x_i^{\alpha}}
   . \]
These are the discrete Hamiltonian equations of motion. In the continuous
limit $\Delta t \to 0$ they reduce to
\[ \dot{x}^{\alpha} (t) = \frac{\partial H}{\partial k_{\alpha}}, \qquad
   \dot{k}_{\alpha} (t) = - \frac{\partial H}{\partial x^{\alpha}}, \]
where $H (k, x) = \hat{r}  (k, x, t)$ is identified as the classical
Hamiltonian (see Section \ref{section:From Quantum to Classical}).

\chapter{Operator Identity: $(p + f)^n = e^{- iF} p^n
e^{iF}$}\label{appendix:operator-identity}

Let $F$ be a smooth function satisfying $F' = f$. The operator $p = - i
\partial_x$ satisfies the commutation relation $[p, f] = - if'$.

\begin{lemma}
  For any smooth function $F$,
  \begin{equation}
    e^{- iF}  \hspace{0.17em} p \hspace{0.17em} e^{iF} = p + F' .
    \label{eq:sim-transform}
  \end{equation}
\end{lemma}

\begin{proof}
  Acting on a test function $\psi$,
  \[ e^{- iF} pe^{iF} \psi = e^{- iF}  (- i \partial_x)  (e^{iF} \psi) = e^{-
     iF}  (F' e^{iF} \psi + e^{iF} (- i \partial_x \psi)) = F' \psi + p \psi .
  \]
  Hence $e^{- iF} pe^{iF} = p + F'$ as operators.
\end{proof}

With $F' = f$, equation \ref{eq:sim-transform} reads $e^{- iF} pe^{iF} = p +
f$. Iterating gives the main result:

\begin{corollary}
  \begin{equation}
    (p + f)^n = e^{- iF}  \hspace{0.17em} p^n  \hspace{0.17em} e^{iF} .
    \label{eq:cor-identity}
  \end{equation}
\end{corollary}

\subsection{Explicit Expansion via Bell Polynomials}

Expanding $p^n e^{iF}$ using Faà di Bruno's formula for $g (F (x)) = e^{iF
(x)}$,
\[ \partial_x^n e^{iF (x)} = e^{iF (x)}  \sum_{k_1 + 2 k_2 + \ldots + nk_n =
   n} \frac{n!}{k_1 ! \hspace{0.17em} k_2 ! \cdots}  \prod_{m = 1}^n \left(
   \frac{iF^{(m)} (x)}{m!} \right)^{k_m} . \]
Multiplying by $(- i)^n$,

\begin{align}
  p^n e^{iF} & = e^{iF}  \sum_{k_1 + 2 k_2 + \ldots + nk_n = n} \frac{n!}{k_1
  ! \hspace{0.17em} k_2 ! \cdots}  (- i)^n  \prod_{m = 1}^n \left(
  \frac{iF^{(m)}}{m!} \right)^{k_m} \nonumber\\
  & = e^{iF}  \sum_{k_1 + 2 k_2 + \ldots + nk_n = n} \frac{n!}{k_1 !
  \hspace{0.17em} k_2 ! \cdots}  \prod_{m = 1}^n \left( \frac{(- i)^m i}{m!}
  \right)^{k_m} \hspace{0.17em} (F^{(m)})^{k_m} . 
\end{align}

Since $(- i)^m i = i^{1 - m}$, the numerical factor simplifies to $i^{1 - m}$
for each $m$-th derivative. Substituting $F' = f$, $F^{(m)} = f^{(m - 1)}$,
and using (\ref{eq:cor-identity}),
\begin{equation}
  (p + f)^n = \sum_{j = 0}^n \binom{n}{j} \sum_{k_1 + 2 k_2 + \ldots + (n - j)
  k_{n - j} = n - j} \frac{(n - j) !}{k_1 ! \hspace{0.17em} k_2 ! \cdots} 
  \prod_{m = 1}^{n - j} \left( \frac{i^{1 - m}}{m!} \right)^{k_m} (f^{(m -
  1)})^{k_m}  \hspace{0.27em} p^j . \label{eq:expansion}
\end{equation}
This can be expressed compactly using the complete Bell polynomials $B_n$:
\begin{equation}
  (p + f)^n = \sum_{j = 0}^n \binom{n}{j} B_{n - j} \left( f, \hspace{0.27em}
  - if', \hspace{0.27em} - f'', \hspace{0.27em} if^{(3)}, \hspace{0.27em}
  f^{(4)}, \hspace{0.27em} \ldots \right)  \hspace{0.27em} p^j .
  \label{eq:bell-expansion}
\end{equation}

\subsection{Examples}

The first few expansions are:

\begin{alignat}{2}
  n=0 & : \qquad 1,\\
  n=1 & : \qquad p+f,\\
  n=2 & : \qquad p\tmrsup{2}+2 f p+ (f\tmrsup{2}-i f\tmrsup{'}),\\
  n=3 & : \qquad p\tmrsup{3}+3 f p\tmrsup{2}+ (3 f\tmrsup{2}-3 i f\tmrsup{'}) 
  p+ (f\tmrsup{3}-3 i f f\tmrsup{'}-f\tmrsup{''}),\\
  n=4 & : \qquad p\tmrsup{4}+4 f p\tmrsup{3}+ (6 f\tmrsup{2}-6 i f\tmrsup{'}) 
  p\tmrsup{2}{\nonumber}\\
  & \qquad\qquad + (4 f\tmrsup{3}-12 i f f\tmrsup{'}-4 f\tmrsup{''}) 
  p{\nonumber}\\
  & \qquad\qquad + (f\tmrsup{4}-6 i f\tmrsup{2} f\tmrsup{'}-4 f
  f\tmrsup{''}+3 f\tmrsup{'}\tmrsup{2}+i f\tmrsup{'''})  .
\end{alignat}

These match the recursive formula derived from $[p, f] = - if'$ given in the
text.

\end{document}
